\section{Conoscenza di dominio}
\label{sec:domain-knowledge}

\subsection{Ordini di vendita e ordini di produzione}
\label{sec:odv-odp}

Un \textbf{ordine di vendita (ODV)} formalizza la richiesta di fornitura di un cliente e ne definisce le condizioni commerciali. Può comprendere più forniture, articolate in distinte posizioni.

Una \textbf{posizione ODV} è la singola voce dell'ordine di vendita che definisce una specifica fornitura attraverso le specifiche del prodotto, la quantità richiesta e le condizioni di consegna. Rappresenta l'unità di domanda a cui viene associata la produzione ed è identificata dalla coppia numero ODV e numero di posizione.

Un \textbf{ordine di produzione (ODP)} formalizza la richiesta interna di realizzare una determinata quantità di prodotto secondo le specifiche produttive assegnate. Costituisce il riferimento per organizzare e registrare le lavorazioni necessarie alla sua realizzazione.

% Domain sources: StrategicPlannerBeanTech/notes/Email Correspondence
% August 24-26 2026.md; notes/Meeting September 29 2026 Decisions.md.
% The August one-to-one statement is superseded by the September description
% and the one-to-many relationships observed in the September snapshot.
Nel contesto considerato, ogni ODP è associato a una e una sola posizione ODV, mentre una stessa posizione ODV può essere associata a uno o più ODP. L'ODV esprime quindi la domanda del cliente; gli ODP ne rappresentano la traduzione in richieste di produzione.

% Example verified directly against StrategicPlannerBeanTech/data/raw/2026-09-24/:
% ordini_di_vendita_260915.csv: Numero_OdV=0001113724, Posizione_OdV=000130;
% ordini_di_produzione_260915.csv: ODV=0001113724, POSIZIONE_ODV=000130.
% Join on both identifiers, preserving leading zeros. Both ODP have material
% 3010148, steel 35MNVS5+CR, dimension 91 mm and requested length 7968 mm.
\begingroup
\setlength{\fboxsep}{10pt}
\par\medskip
\noindent\colorbox{ReportBlue!7}{%
\begin{minipage}{\dimexpr\linewidth-2\fboxsep\relax}
\small
\textbf{Esempio dal dataset: una posizione ODV associata a due ODP}

Nell'estrazione ricevuta il 24 settembre 2026, la posizione \texttt{000130} dell'ODV \texttt{0001113724} riguarda il materiale \texttt{3010148}, con data di consegna richiesta pari al 2 gennaio 2026. Alla posizione sono associati i seguenti ordini di produzione:

\begin{center}
\begin{tabular}{ll}
\toprule
\textbf{ODP} & \textbf{Colata} \\
\midrule
\texttt{000010874375} & \texttt{298932} \\
\texttt{000010880354} & \texttt{299558} \\
\bottomrule
\end{tabular}
\end{center}

\end{minipage}%
}
\par\medskip
\endgroup

\subsection{Macrofasi del processo produttivo}
\label{sec:production-phases}

Per realizzare il prodotto richiesto da un ODP, il materiale attraversa un percorso di lavorazioni e controlli, fino al suo approntamento. Le registrazioni che documentano quanto avvenuto durante questo percorso sono dette \textbf{consuntivi di produzione}. Esse associano all'ODP le attività e gli stati registrati, con le relative informazioni temporali e quantitative, e permettono di ricostruirne la storia produttiva.

Una \textbf{macrofase} è un passaggio di questo percorso, descritto a un livello che riassume il dettaglio delle singole operazioni. Può riunire più lavorazioni in un'unica fase, identificare un controllo oppure rappresentare uno stato del materiale, come l'attesa di un esito o il completamento della produzione. Il termine indica quindi il ruolo del passaggio nel processo; non implica che esso corrisponda a una sola macchina.

Nei consuntivi, ogni registrazione riferita a una macrofase ne descrive una specifica occorrenza per un ODP. Una stessa macrofase può comparire più volte nel percorso dello stesso ordine. La successione dei passaggi registrati descrive il \textbf{ciclo produttivo osservato} dell'ODP.

% Inventory checked directly against all records in
% StrategicPlannerBeanTech/data/raw/2026-09-24/Consuntivi_260915.csv,
% field MACHINE_AGGREGATION_CODE: 14 distinct values, 35893 records.
% This is the raw September snapshot, not the August inventory or the
% physically filtered September baseline.
% Domain sources:
% - data/raw/2026-09-24/ABS_Flussi_Luna.pptx, slides 1--8;
% - notes/Email Correspondence August 24-26 2026.md;
% - notes/Meeting August 27 2026 Summary.md;
% - notes/Meeting September 29 2026 Decisions.md.
% Earlier interpretations and unresolved mappings were cross-checked with
% docs/entity_model_specification.md and
% notebooks/process understanding/03 - phase ontology.ipynb.
Nei consuntivi dell'estrazione ricevuta il 24 settembre 2026 sono presenti 14 etichette distinte. Il loro significato è ricostruito a partire dal documento aziendale \emph{ABS\_Flussi\_Luna} e dalle note di progetto. 

\subsubsection{Processo iniziale aggregato}

\begin{description}
\item[\texttt{TRENOLUNA}.] Comprende le lavorazioni in continuo del Treno Luna, dal riscaldo al controllo qualità in linea. Il diagramma aziendale distingue al suo interno la laminazione, il percorso di raffreddamento o trattamento e il controllo del materiale. Nei consuntivi queste attività sono riunite in un'unica macrofase, precedente alle eventuali lavorazioni successive e ai percorsi di recupero del materiale non conforme.
\end{description}

\subsubsection{Lavorazioni e controlli}

\begin{description}
\item[\texttt{BANCHI}.] Indica il controllo del materiale ai banchi, ricondotto nel diagramma aziendale alla voce banco/metalloscopio. Serve a verificare il materiale e a indirizzare gli eventuali interventi di recupero dopo il controllo in linea.

\item[\texttt{TAG}.] Identifica le operazioni di taglio, rappresentate nella documentazione dalla voce \emph{taglio seghetti}. Può comparire nei percorsi di rilavorazione del materiale.

\item[\texttt{MOL}.] Indica la molatura, ovvero una lavorazione meccanica del materiale e compare negli schemi di recupero insieme al taglio e alla raddrizzatura.

\item[\texttt{RADDR}.] Indica la raddrizzatura, ossia la lavorazione destinata a correggere la rettilineità del materiale. Nel diagramma compare sia nei percorsi di recupero sia dopo i trattamenti termici.

\item[\texttt{CIRCO}.] È interpretata come il controllo mediante \emph{CIRCOGRAPH}, descritto nel PowerPoint come controllo qualità finale per le rilavorazioni. La corrispondenza tra il codice abbreviato del dataset e questa voce del diagramma resta da confermare.

\item[\texttt{CONF}.] È interpretata come confezionamento, coerentemente con la voce \emph{Confez} presente nei diagrammi in prossimità dell'uscita del materiale. La documentazione disponibile non esplicita la corrispondenza con il codice \texttt{CONF}, che resta da confermare.
\end{description}

Il ruolo di una lavorazione dipende dal percorso del prodotto: la presenza di una macrofase non basta a qualificarla come rilavorazione. La raddrizzatura, per esempio, può essere un passaggio previsto dopo un trattamento termico oppure un intervento di recupero.

\subsubsection{Trattamenti}

\begin{description}
\item[\texttt{TRATT\_LUNA}.] Identifica una macrofase di trattamento riferita a Luna. Il documento aziendale descrive trattamenti termici a induzione, ma non fornisce una mappatura esplicita tra queste attività e il codice del dataset. Il tipo di trattamento e la risorsa utilizzata non possono quindi essere dedotti dalla sola etichetta.

\item[\texttt{TRATT\_LFC}.] Identifica un'ulteriore macrofase di trattamento, distinta nei dati da \texttt{TRATT\_LUNA}. Il diagramma comprende anche trattamenti termici esterni a Luna, ma non ne stabilisce la corrispondenza con questo codice. Il significato esteso di \texttt{LFC} e la natura interna o esterna del trattamento restano da chiarire.
\end{description}

\subsubsection{Controlli di rilascio e stati di avanzamento}

\begin{description}
\item[\texttt{HPT}.] Rappresenta un punto di attesa per il controllo qualità fuori linea, svolto in laboratorio dopo un trattamento termico. Il materiale attende l'esito dei controlli prima di proseguire. Nei consuntivi questa macrofase ha durata registrata nulla: l'attesa dei risultati può comunque separare temporalmente le attività adiacenti.

\item[\texttt{ANTIMIX}.] Nel documento aziendale indica il controllo volto a rilevare mescolamenti qualitativi del materiale. Nelle indicazioni di progetto viene trattata come un punto di controllo o attesa, senza un vincolo di capacità proprio, anche quando il consuntivo riporta una durata positiva. Occorre quindi distinguere il significato fisico del controllo dalla sua rappresentazione nel modello di pianificazione.

\item[\texttt{PRONTO}.] È lo stato terminale che indica l'approntamento del materiale al termine del percorso registrato. Identifica il completamento produttivo rappresentato nei consuntivi, senza attestare di per sé l'avvenuta spedizione al cliente.

\item[\texttt{PRONTO\_QS}.] È una variante dello stato terminale di approntamento. Il significato specifico del suffisso \texttt{QS} non è definito nelle fonti disponibili e viene mantenuto senza attribuirgli una lavorazione aggiuntiva.

\item[\texttt{PRONTO\_QWR}.] È un'ulteriore variante dello stato terminale di approntamento. Anche il significato del suffisso \texttt{QWR} richiede un chiarimento aziendale; nei consuntivi resta distinto dalle altre due etichette terminali.
\end{description}

% Real example: September Consuntivi_260915.csv, ODP=000010858580,
% ODV=0001117836, POSIZIONE_ODV=000020, source CSV lines 22--25.
% Operation order is the source-row order.
% Recorded route: TRENOLUNA > BANCHI > ANTIMIX > PRONTO.
% ANTIMIX has DURATA_TOT_MIN=18; this is recorded duration, not waiting time.
\begingroup
\setlength{\fboxsep}{10pt}
\par\medskip
\noindent\colorbox{ReportBlue!7}{%
\begin{minipage}{\dimexpr\linewidth-2\fboxsep\relax}
\small
\textbf{Esempio dal dataset: successione di macrofasi}

L'ODP \texttt{000010858580} presenta nei consuntivi il percorso
\[
\texttt{TRENOLUNA}\;\rightarrow\;\texttt{BANCHI}\;\rightarrow\;
\texttt{ANTIMIX}\;\rightarrow\;\texttt{PRONTO}.
\]
Al processo iniziale aggregato seguono il controllo ai banchi, il controllo antimix e lo stato finale di approntamento. Per \texttt{ANTIMIX} sono registrati 18 minuti di durata, mentre \texttt{PRONTO} ha durata nulla.
\end{minipage}%
}
\par\medskip
\endgroup
